\subsection{Summary of work}

The main goal of this experiment is to study the robustness of different NMF methods in noisy environments. We used the $L_2$ NMF method and the $L_{2,1}$ norm based NMF method to study the performance of the datasets ORL and YaleB in the salt and pepper noise environment. We used the indicators Relative Reconstruction Error(RRE) and Accuracy/Normalized Mutual Information(Acc/NMI) to evaluate our results.

\subsection{Key findings}

From the experimental results, several trends are evident. 
First, both $L_2$ and $L_{2,1}$ NMF algorithms converge stably under all tested noise levels, but their performance differs across datasets and noise intensities. 
On the smaller ORL dataset, the $L_{2,1}$-norm NMF generally achieves lower reconstruction error (RRE) and slightly better clustering accuracy under low noise conditions, confirming its robustness to impulsive salt-and-pepper corruption. 
On the larger YaleB dataset, the two algorithms perform comparably in terms of RRE, while $L_{2,1}$ NMF shows slightly higher clustering accuracy when the noise ratio is moderate. 
As the noise proportion increases, both methods show clear degradation in accuracy and normalized mutual information (NMI), though $L_{2,1}$ tends to maintain more stable reconstruction performance. 
Overall, noise level has a strong and consistent impact on both algorithms, especially on the larger YaleB dataset.

\subsection{Interpretation and Insights}

The comparative results highlight complementary strengths of the two algorithms. 
The $L_{2,1}$-norm NMF demonstrates stronger robustness to sparse, high-intensity noise, as its row-sparse penalty suppresses the influence of corrupted samples, explaining its lower reconstruction errors on ORL. 
However, the $L_2$-norm NMF tends to preserve more global structure, which benefits clustering on large-scale datasets such as YaleB when the noise distribution is relatively uniform. 
These observations are consistent with previous findings \cite{Zeng2023RobustnessThreeNMF, Shen2022AdaptiveWeightedNMF, Liu2022RLNMF_AG}, indicating that robust NMF variants are advantageous under strong impulsive noise but not necessarily superior for clean or weakly corrupted data. 
In summary, $L_{2,1}$ is preferable when robustness is the main goal, whereas $L_2$ remains effective for large, structured datasets where noise is moderate and the focus is clustering performance.

\subsection{Limitations and Future Work}

We encountered many problems in our work, including:

\begin{itemize}
    \item The computational complexity is high, the operation time is long, and the convergence speed is slow.
    \item Few types of noise were tested. 
    \item The NMF method is highly sensitive to parameters.
\end{itemize}

\bigskip

We consider the following methods to improve the next step of the experiment:

\begin{itemize}
    \item Combined with the $L_1$ norm, a hybrid model is formed to evaluate its robustness.
    \item Introducing sparsity or graph regularization
    \item Use larger-scale or online updated NMF algorithms
\end{itemize}

\subsection{Closing Statement}

In summary, this paper verifies the robustness advantage of $L_{2,1}$ norm in non-negative matrix factorization through the study of $L_2$ NMF method and $L_{2,1}$ NMF method. The $L_{2,1}$ method can effectively maintain data structure and clustering performance in a high-noise environment. It provides a reliable theoretical and practical basis for real-world projects such as face recognition.